% ----- CHAUPAI 19 -----

\newpage

\subsection*{\deva{एकोनविंशतितम चौपाई} \(Nineteenth Chaupai\)}
\addcontentsline{toc}{subsection}{\deva{एकोनविंशतितम चौपाई} \(Nineteenth Chaupai\) — \deva{प्रभु मुद्रिका...}}

\begin{minipage}[t]{0.55\textwidth}
\vspace{0pt}

{\large \linespread{1.5}\selectfont
\deva{प्रभु मुद्रिका मेलि मुख माहीं।}\\
\deva{जलधि लांघि गये अचरज नाहीं॥}
\par}

\vspace{0.5em}

{\large \linespread{1.5}\selectfont
\telu{ప్రభు ముద్రికా మేలి ముఖ మాహీఁ।}\\
\telu{జలధి లాంఘి గయే అచరజ నాహీఁ॥}
\par}

\vspace{0.5em}

{\large \linespread{1.3}\selectfont
\textit{prabhu mudrikā meli mukha māhīṃ |}\\
\textit{jaladhi lāṃghi gaye acaraja nāhīṃ ||}
\par}
\end{minipage}%
\hfill
\begin{minipage}[t]{0.42\textwidth}
\vspace{0pt}
\begin{tcolorbox}[anchorbox]
\textbf{Focus Under Pressure:} He carried the absolute weight of a massive responsibility (Rama's signet ring) across a terrifying, treacherous ocean without dropping it. Taking on great responsibility in education or life requires unshakeable focus and the determination not to let distractions sink your goals.
\vspace{0.5em}\newline He carried an enormous responsibility across a treacherous, demon-patrolled ocean without ever losing focus.\newline\null\hfill\href{https://www.valmiki.iitk.ac.in/sloka?field_kanda_tid=5&language=dv&field_sarga_value=36}{\textit{Sundara Kanda, 36:2}}
\end{tcolorbox}
\end{minipage}

\vspace{0.5em}
\subsubsection*{\textbf{Visual Rhythm Map:}}
\begin{table}[h!]
\centering
\renewcommand{\arraystretch}{1.2}
\setlength{\tabcolsep}{0pt}
\begin{tabularx}{\textwidth}{|*{16}{>{\centering\arraybackslash\hsize=1.0\hsize}X|}}
\hline
\multicolumn{2}{|>{\centering\arraybackslash\hsize=2.0\hsize}X|}{\textbf{\laghu{प्र}\laghu{भु}} \par (2)} &
\multicolumn{5}{>{\centering\arraybackslash\hsize=5.0\hsize}X|}{\textbf{\guru{मु}\laghu{द्रि}\guru{का}} \par (5)} &
\multicolumn{3}{>{\centering\arraybackslash\hsize=3.0\hsize}X|}{\textbf{\guru{मे}\laghu{लि}} \par (3)} &
\multicolumn{2}{>{\centering\arraybackslash\hsize=2.0\hsize}X|}{\textbf{\laghu{मु}\laghu{ख}} \par (2)} &
\multicolumn{4}{>{\centering\arraybackslash\hsize=4.0\hsize}X|}{\textbf{\guru{मा}\guru{हीं}} \par (4)} \\
\hline
\end{tabularx}
\vspace{-1.5em}
\end{table}

\begin{table}[h!]
\centering
\renewcommand{\arraystretch}{1.2}
\noindent\makebox[\textwidth][l]{%
  {%
  \setlength{\tabcolsep}{0pt}%
  \begin{tabularx}{1.0625\textwidth}{|*{17}{>{\centering\arraybackslash\hsize=1.0\hsize}X|}}
  \hline
  \multicolumn{3}{|>{\centering\arraybackslash\hsize=3.0\hsize}X|}{\textbf{\laghu{ज}\laghu{ल}\laghu{धि}} \par (3)} &
  \multicolumn{3}{>{\centering\arraybackslash\hsize=3.0\hsize}X|}{\textbf{\guru{लां}\laghu{घि}} \par (3)} &
  \multicolumn{3}{>{\centering\arraybackslash\hsize=3.0\hsize}X|}{\textbf{\laghu{ग}\guru{ये}} \par (3)} &
  \multicolumn{4}{>{\centering\arraybackslash\hsize=4.0\hsize}X|}{\textbf{\laghu{अ}\laghu{च}\laghu{र}\laghu{ज}} \par (4)} &
  \multicolumn{4}{>{\centering\arraybackslash\hsize=4.0\hsize}X|}{\textbf{\guru{ना}\guru{हीं}} \par (4)} \\
  \hline
  \end{tabularx}%
  }%
  \hspace{0.01\textwidth}%
  \begin{minipage}{0.1\textwidth}
  \vspace{-0.5em}
  \centering
  {\Huge \textbf{17}}
  \end{minipage}%
}
\vspace{-1.5em}
\end{table}

\vspace{0.5em}
\textbf{ \deva{सरल-संस्कृतम्}:}\\
 \deva{प्रभोः अङ्गुलीयकं (मुद्रिकाम्) मुखे निधाय }\\
 \deva{भवान् समुद्रं (जलधिं) उल्लङ्घितवान्, अत्र किञ्चित् आश्चर्यं नास्ति।}

Keeping the Lord's ring in your mouth you leaped across the vast ocean, 
and there is no surprise in this!

\vspace{1em}
\begin{tcolorbox}[poetrybox]
\textbf{The 17-Beat Squeeze \& Singing Flow:}\\
Line 2 mathematically counts to 17 beats because \deva{गये} is fully 3 beats native. But because this poem is \textit{Bhava Pradhan} (driven by devotion), strict meter bends to the singer's emotion. Singers effortlessly compress \deva{गये} into 2 quick beats (\textit{ga-e}). The 16-beat rule is a flexible guide, and these anomalies are part of the lively, organic beauty of the Awadhi oral tradition!
\end{tcolorbox}
\newpage
\textbf{\deva{प्रतिपदार्थः} (Meaning of each word):}
\begin{table}[h!]
\centering
\renewcommand{\arraystretch}{1.2}
\begin{tabularx}{\textwidth}{@{} l l X X @{}}
\toprule
\textbf{अवधी} & \textbf{संस्कृतम्} & \textbf{English} & \textbf{Grammar \& Notes} \\
\midrule
\deva{प्रभु} / \telu{ప్రభు}  &  \deva{प्रभोः / रामस्य}  &  Lord's  &   \\
\deva{मुद्रिका} / \telu{ముద్రికా}  &  \deva{मुद्रिका / अङ्गुलीयकम्}  &  Ring  &   \\
\deva{मेलि} / \telu{మేలి}  &  \deva{निधाय}  &  Having kept  & Absolutive `-i` ending \\
\deva{मुख माहीं} / \telu{ముఖ మాహీఁ}  &  \deva{मुख-मध्ये}  &  Inside the mouth  &   \\
\deva{जलधि} / \telu{జలధి}  &  \deva{जलधिः (समुद्रः)}  &  Ocean  &   \\
\deva{लांघि} / \telu{లాంఘి}  &  \deva{उल्लङ्घ्य}  &  Having crossed/leapt  &   \\
\deva{गये} / \telu{గయే}  &  \deva{गतवान्}  &  Went  &   \\
\deva{अचरज} / \telu{అచరజ}  &  \deva{आश्चर्यम्}  &  Surprise  & Swara-bhakti \\ 
\deva{नाहीं} / \telu{నాహీఁ}  &  \deva{नास्ति}  &  Is not  &   \\
\bottomrule
\end{tabularx}
\end{table}



\newpage
